% !TEX root = ../../tir_monograph_v12.tex
\chapter{Formal Proofs, Algebraic Firewalls and Legacy Repairs}
\label{app:v12_formal_proofs}

\paragraph{Purpose and provenance.}
This appendix consolidates long algebra that supports Chapters~4--10 and the
diagnostic no-go surfaces used later in the monograph.  Historical appendix
material is preserved through explicit source pointers rather than copied
verbatim.  GREMLIN/HOUND is used here as a bounded adversarial audit layer:
candidate inconsistencies are promoted into this appendix only after an exact
algebraic or deterministic check.

The principal legacy sources are
\path{TIR/monograph/appendices/appB_collatz_tables.tex},
\path{TIR/monograph/appendices/appC_su3_primer.tex},
\path{TIR/monograph/appendices/appD_poincare_disk.tex},
\path{TIR/monograph/appendices/appE_tetrahedron_algebra.tex},
\path{TIR/monograph/appendices/appJ_collatz_quarter_power_scaling.tex},
\path{TIR/monograph/appendices/appK_sector_holonomy_release.tex}, and
\path{TIR/monograph/appendices/appL_up_sector_common_baseline_no_go.tex}.

\section{Finite Collatz statements and the structural labels}

For positive integers define
\[
C(n)=
\begin{cases}
3n+1,&n\ \hbox{odd},\\
n/2,&n\ \hbox{even}.
\end{cases}
\]
The seed required by the canonical v12 label is finite and directly checkable:
\[
3\to10\to5\to16\to8\to4\to2\to1.
\]
There are seven steps before the first occurrence of \(1\), giving the typed
structural label
\[
\boxed{L_3=7.}
\]
The twin-prime source gives
\[
\boxed{L_4=5-3=2,\qquad L_5=5.}
\]
These finite statements are the complete arithmetic input required by the
v12 label layer.  Their publication status is inherited from
Chapter~\ref{ch:v12_discrete_labels}.

A second finite trajectory, historically used to motivate the quark-prime
candidate set, is
\[
7\to22\to11\to34\to17\to52\to26\to13\to40\to20\to10\to5
\to16\to8\to4\to2\to1.
\]
The primes encountered are
\[
7,11,17,13,5,3,
\]
whose sorted set is
\[
\mathcal P_Q=\{3,5,7,11,13,17\}.
\]
This is retained as a finite set-generation observation.  The assignment to
ordered quark slots is owned by the explicit typed map of
Chapter~\ref{ch:v12_discrete_labels}; the exhaustive \(6!=720\) permutation
audit leaves two arithmetic survivors before that ordering rule and one after
it.

\section{Corrected \(SU(3)\) representation algebra}

The three-flavour carrier is \(V_F\cong\mathbb C^3\), with mixing algebra
\(\mathfrak{su}(3)\).  Its real dimension is
\[
\dim_{\mathbb R}\mathfrak{su}(3)=3^2-1=8.
\]
For an \(SU(3)\) irreducible representation with Dynkin labels \((p,q)\),
\[
\boxed{
\dim(p,q)=\frac{(p+1)(q+1)(p+q+2)}{2}.
}
\]
Consequently,
\[
\dim(1,0)=3,\qquad
\dim(0,1)=3,\qquad
\dim(1,1)=8,\qquad
\dim(3,0)=10,\qquad
\dim(2,2)=27.
\]
The mesonic flavour nonet belongs to the reducible decomposition
\[
\boxed{
3\otimes\bar3=8\oplus1,
}
\]
rather than to the \((2,2)\) irrep.  The historical primer row that paired
\((2,2)\) with ``meson nonet'' is therefore retained as a repaired
representation-label mismatch.
\vTwelveStatus{D}{RETROSPECTIVE}{PASS}

For the canonical \(\kappa\) normalization, only the exact carrier facts
\[
N_F=3,\qquad \dim\mathfrak{su}(3)=8,\qquad N_{\rm mix}=3\times8=24
\]
are required.  The full normalization derivation remains owned by
Chapter~\ref{ch:v12_kappa_normalization} and
\path{TIR/foundations/TIR_KAPPA_FLAVOUR_MIXING_NORMALIZATION_V0_1.md}.

\section{Poincar\'e disk normalization repair}

For the unit disk \(\mathbb D=\{z\in\mathbb C:|z|<1\}\) with metric
\[
ds^2=\frac{4\,|dz|^2}{(1-|z|^2)^2},
\]
the curvature is \(K=-1\).  Define the disk cross-ratio modulus
\[
\chi(z_1,z_2)=
\left|
\frac{z_1-z_2}{1-\bar z_1z_2}
\right|.
\]
The geodesic distance compatible with the displayed metric is
\[
\boxed{
d_{\mathbb D}(z_1,z_2)
=2\,\operatorname{artanh}\chi(z_1,z_2).
}
\]
For \(z_1=0\) this reduces to
\[
d_{\mathbb D}(0,r)=\int_0^r\frac{2\,du}{1-u^2}
=2\,\operatorname{artanh}r,
\]
which fixes the normalization directly.  The historical appendix printed the
same expression without the factor \(2\); v12 records that line as a metric
normalization repair.
\vTwelveStatus{D}{RETROSPECTIVE}{PASS}

For a geodesic triangle with angles \(\alpha,\beta,\gamma\),
\[
A_{\mathbb D}=\pi-(\alpha+\beta+\gamma).
\]
A one-angle identification therefore supplies only partial cell data.  v12
assigns any proposed information-area equality to an explicit bridge theorem
with the complete triangle, curvature convention and area normalization stated
before comparison.

\section{Regular tetrahedral Gram algebra}

Let \(\mathbf n_a\in\mathbb R^3\), \(a=1,\ldots,4\), be unit vectors satisfying
\[
\mathbf n_a\cdot\mathbf n_b=-\frac13\qquad(a\ne b).
\]
Their Gram matrix is
\[
G_{ab}=
\begin{cases}
1,&a=b,\\
-1/3,&a\ne b.
\end{cases}
\]
Equivalently,
\[
G=\frac43I_4-\frac13J_4,
\]
where \(J_4\) is the all-ones matrix.  Since \(J_4\) has eigenvalues \(4,0,0,0\),
\[
\operatorname{spec}(G)=\left\{0,\frac43,\frac43,\frac43\right\},
\qquad
\operatorname{rank}G=3.
\]
The null eigenvector is \((1,1,1,1)^T\), hence
\[
\sum_{a=1}^4\mathbf n_a=0.
\]
Every squared edge length is
\[
|\mathbf n_a-\mathbf n_b|^2
=2-2\mathbf n_a\cdot\mathbf n_b
=\frac83,
\]
so the four points form a regular tetrahedron in the three-real-dimensional
carrier.  Permuting the four vertices preserves the Gram matrix, giving
\[
\operatorname{Aut}(\Delta^3)\cong S_4,
\qquad |S_4|=24.
\]
This is the finite-symmetry crosscheck used after the independent
flavour-mixing count \(3\times8=24\).

\section{Legacy operator-tetrahedron firewall}

The historical operator-labelled tetrahedron introduced six sector-specific
edge numbers and then wrote
\[
\cos\theta_{ij}
=
\frac{e_{ik}^2+e_{jk}^2-e_{ij}^2}{2e_{ik}e_{jk}}.
\]
That expression is the planar law of cosines for the angle of a triangle built
from three edge lengths.  A tetrahedral dihedral angle instead requires the
normals of two faces, or an equivalent Cayley--Menger/Gram construction using
the full tetrahedral metric data.  v12 therefore assigns the historical
``dihedral'' line to the diagnostic ledger rather than to the regular
tetrahedral closure theorem.
\vTwelveStatus{D}{RETROSPECTIVE}{QUARANTINED}

The same historical appendix states a volume proportionality to \(\kappa^3\)
without an independently derived Gram determinant and scale map.  v12 assigns
volume/PMNS coupling to the sector-model provenance of Chapters~12--13, while
the geometric tetrahedron theorem is owned by Chapter~6.

\section{Accelerated Collatz quarter-power arithmetic}

For odd \(n\), define the accelerated map
\[
T(n)=\frac{3n+1}{2^{a(n)}},
\qquad
a(n)=\nu_2(3n+1).
\]
Under the standard uniform odd-residue calculation,
\[
\Pr(a=k)=2^{-k},\qquad k\ge1,
\]
so
\[
\mathbb E[a]=\sum_{k\ge1}k2^{-k}=2.
\]
The associated asymptotic geometric-mean multiplier is then
\[
\boxed{
\rho_C
=\exp\!\left(\ln3-\mathbb E[a]\ln2\right)
=\frac34.
}
\]
This arithmetic result is distinct from the particle-model bridge that uses
\(3/4\) as an exponent.  The archived one-anchor charged-fermion trace keeps its
retrospective physical-spectrum failure, while the arithmetic multiplier
remains a technical result.

\section{Common-baseline translation no-go}

Let \(r_g\) be logarithmic residuals for a fixed three-channel up-sector trace,
and apply one common additive baseline \(B\):
\[
r'_g=r_g+B.
\]
For every pair \(g,h\),
\[
r'_g-r'_h=r_g-r_h,
\]
so the residual spread is translation invariant.

\paragraph{Proposition.}
There exists a common \(B\) satisfying
\[
|r_g+B|\le\varepsilon
\quad\text{for every }g
\]
exactly when
\[
\max_g r_g-\min_g r_g\le2\varepsilon.
\]
The smallest achievable uniform error is
\[
\boxed{
\varepsilon_{\min}
=\frac{\max r-\min r}{2}.
}
\]

\paragraph{Proof.}
Each channel requires
\[
B\in[-r_g-\varepsilon,-r_g+\varepsilon].
\]
All intervals intersect exactly when the largest lower endpoint is no greater
than the smallest upper endpoint, which is equivalent to the stated spread
condition.  Centering the two extreme residuals around zero attains half the
spread.

For the archived v10.2 residual vector
\[
(r_u,r_c,r_t)
=(3.1857238256,\,0.0818052322,\,0.2293790515),
\]
the spread is
\[
\Delta r=3.1039185934
\]
and therefore
\[
\varepsilon_{\min}=1.5519592967,
\qquad
e^{\varepsilon_{\min}}=4.7207104\ldots .
\]
The common-baseline architecture is thus excluded for the fixed relative trace
at any target envelope narrower than this minimax bound.
\vTwelveStatus{D}{RETROSPECTIVE}{PASS}

\section{Appendix-A audit ownership}

The deterministic appendix integration gate is
\path{TIR/validation/tir_v12_appendix_integration_audit_v0_1.py}.
Its machine-readable receipt is
\path{TIR/validation/TIR_V12_APPENDIX_INTEGRATION_AUDIT_V0_1.json}.
The gate checks the corrected \(SU(3)\) representation identity, the factor-two
Poincar\'e distance normalization, the regular tetrahedral Gram spectrum, the
legacy dihedral quarantine marker and the common-baseline minimax identity.
